
suite les immersions @@M0@@ (@@M1@@) sont régulières, et restent des immersions régulières @@M2@@ par changement de base @@M3@@. Notons que en vertu de la suite exacte d'homotopie (IX 1.1), @@M4@@ provient d'un élément @@M5@@, donc l'élément @@M6@@ de @@M7@@ provient de l'élément @@M8@@ de @@M9@@, et par suite son image dans @@M10@@ par (7.3.1) n'est autre que
@@M11@@
Bien entendu, pour une section @@M12@@ de @@M13@@ sur @@M14@@, on pose
@@M15@@
est défini comme plus haut, et peut être considéré comme déduit de @@M16@@ par le changement de base plat @@M17@@; c'est donc une immersion régulière, à fortiori un morphisme parfait, ce qui donne bien un sens à l'expression @@M18@@. Pour la même raison l'expression @@M19@@ est définie, et par transitivité des images inverses (IV 2.12) dans le diagramme commutatif de morphismes parfaits suivant :
@@M20@@
on trouve que le premier membre de @@M21@@ est égal à
@@M22@@
Comme @@M23@@ est algébriquement clos, les composantes irréductibles réduites de @@M24@@ sont géométriquement intègres sur @@M25@@ et leurs points d'intersections rationnels sur @@M26@@. On en conclut aisément que l'expression @@M27@@ est algébriquement équivalente à zéro dans @@M28@@, en utilisant la connexité
